\section{Introduction}
\label{sec:introduction}

Single-cell RNA sequencing (scRNA-seq) enables transcriptome-wide measurements
at cellular resolution, but converting them into biologically meaningful cell
populations remains difficult.  Expression matrices are high-dimensional,
sparse, and affected by sequencing depth, technical noise, transient states,
and class imbalance.  Unsupervised clustering therefore depends critically on
a representation that suppresses nuisance variation without erasing signals of
rare types or continuous transitions~\cite{kiselev2019challenges}.  Deep
generative and autoencoding models learn nonlinear structure from expression
profiles~\cite{lopez2018scvi,tian2019scdeepcluster,tian2021scdcc}, yet their
objectives must still determine which variations should be reconstructed,
ignored, or treated as local biological organization.

Masked representation learning offers a direct self-supervised objective.  A
masked autoencoder corrupts part of an input and infers the replaced content
from its context, a principle successful in other modalities~\cite{he2022mae}.
scMAE adapted it to scRNA-seq by jointly reconstructing gene expression and
predicting the corruption mask~\cite{fang2024scmae}.  Its replacement noise
uses within-cell gene context effectively, but it does not explicitly exploit
relationships among cells.  Local expression variation could provide a more
structured perturbation, asking the encoder which anchor-cell features remain
stable among phenotypically related cells.

This opportunity is also a hazard.  A nearest-neighbor graph estimated from
sparse expression can cross cell types, transitional states, technical batches,
or rare-population boundaries.  Uniform neighborhood mixing can therefore move
an anchor away from its biological support and encourage excessive smoothing.
The central question is not merely whether neighbors should be used, but how a
local perturbation should be constructed, bounded, and weighted when the graph
itself is uncertain.

We introduce \emph{single-cell Vicinal Corruption with Anchor Recovery}
(scVICAR), a graph-vicinal extension of masked single-cell representation
learning.  Given a row-stochastic sampled neighborhood operator $\widehat P$
and a diagonal matrix $G$ of cell-wise perturbation budgets, scVICAR constructs
\begin{equation}
 X^v=(I-G)X+G\widehat PX=X-G(I-\widehat P)X.
 \label{eq:intro-view}
\end{equation}
The view is then subjected to the original gene-replacement corruption, while
the reconstruction target remains the unmodified anchor cell.  This
\emph{anchor-recovery} objective differs from conventional mixup, which
interpolates inputs and prediction targets together~\cite{zhang2018mixup}.  The
clean scMAE branch is retained, so vicinal recovery is an auxiliary regularizer
rather than a replacement objective.

scVICAR has two matched variants.  scVICAR-F uses a fixed perturbation budget.
scVICAR-T derives topology-informed affinity from cosine similarity,
mutual-KNN relations, and shared-nearest-neighbor overlap, and assigns each cell
an analytic perturbation budget from local graph coherence.  The topology score
is not a calibrated probability of biological correctness.  It defines a
bounded trust region in which neighborhood information can influence the
corrupted view.  We therefore frame scVICAR-T as a robustness extension of
scVICAR-F under graph contamination, not as a method expected to dominate the
fixed variant on every clean dataset.

The formulation admits several precise interpretations.  Every vicinal view
lies in the convex hull of its anchor and sampled neighbors, and its displacement
is bounded by the cell-wise gate and neighborhood radius.  Equation
\eqref{eq:intro-view} is one step of nonuniform graph diffusion.  A conditional
local expansion of the anchor-recovery loss further gives a directional
curvature-regularization interpretation.  These properties connect local
corruption to bounded denoising without assuming that every graph edge joins
cells of the same type.

We evaluate scVICAR with a frozen, matched-backbone protocol.  NoMix,
scVICAR-F, edge-only, gate-only, topology-adaptive, and random-mixing controls
share preprocessing, encoder, optimizer, training duration, and seed schedule.
Six confirmatory datasets were excluded from the historical NeighborMix and RG
parameter-development stage.  Dataset, rather than seed run or cell, is the
independent statistical unit.  Controlled bad-edge injection tests whether
topology adaptation degrades more gracefully as graph quality falls.  In the
clean primary protocol, labels are used only for predefined filtering,
evaluation, and post-hoc diagnostics; they never enter training, graph
construction, or parameter selection.  The controlled stress intervention uses
labels solely to inject prespecified cross-class edges and is never used to fit
or select a clean model.

Clustering agreement alone does not establish biological utility.  We therefore
complement clustering scores with marker recovery, marker-overlap annotation,
rare-class diagnostics, and frozen-embedding low-label probes, following the
task-diverse evaluation motivation of scCluBench~\cite{xu2026scclubench}.  The
low-label analyses are explicitly within-dataset and transductive, not claims of
cross-dataset transfer.

Our contributions are fourfold:
\begin{itemize}
 \item We formulate graph-vicinal anchor recovery, which uses a bounded local
 view while preserving the original cell as the self-supervised target.
 \item We separate topology-informed edge affinity from a cell-wise perturbation
 budget, with robustness to graph contamination as the adaptive variant's
 primary hypothesis.
 \item We derive convex-hull, perturbation-radius, nonuniform-diffusion, and
 contamination bounds together with a conditional local-risk interpretation.
 \item We establish a matched, confirmatory evaluation spanning clustering,
 controlled graph corruption, and downstream biological utility without
 claiming universal superiority of scVICAR-T over scVICAR-F.
\end{itemize}
